%%%%%%%%%%%%%%%%%%%%%%%%%%%%%%%%%%%%%%%%%%%%%%%%%%%%%%%%%%%%%%%%%%%%%%%%
\section{Meanders with bounded labels}\label{sec:band}
%%%%%%%%%%%%%%%%%%%%%%%%%%%%%%%%%%%%%%%%%%%%%%%%%%%%%%%%%%%%%%%%%%%%%%%%

For all meanders the exponent is out of reach. For one natural subfamily both the growth rate and the
exponent can be found exactly, and the answer shows a mechanism that is believed to produce a power
law in general: a walk confined to a wedge.

\subsection{Labels}

In this section $m$ is the order of a meander. Give each crossing $i<2m$ an integer label $\varphi_i$
so that along every arch above the river the label drops by one from left to right, and along every
arch below the river it rises by one. Precisely: $\varphi_1=0$, and $\varphi_y=\varphi_x-1$ for every
upper arch $\{x,y\}$ with $x<y<2m$, and $\varphi_y=\varphi_x+1$ for every lower arch $\{x,y\}$ with
$x<y<2m$; the two arches at the last crossing $2m$ are left out. Removing them leaves the road as a
single path through crossings $1,\ldots,2m-1$, so these rules determine the labels uniquely.
Intuitively, $\varphi_i$ counts the half turns of the road on its way from crossing $1$ to crossing
$i$, with opposite signs above and below the river. For a meander, labels of neighboring crossings
along the river differ by one. Conversely, a pair of arch systems is a meander if and only if it has a
labeling that obeys these rules and in which neighboring labels differ by one. Paper~B explains these
labels~\cite{PaperB}. In the running example of Figure~\ref{fig:running} the labels of crossings $1$
to $5$ are $0,-1,0,1,0$.

Let $M^{[-1,2]}_m$ be the number of meanders of order $m$ whose labels all lie in $\{-1,0,1,2\}$. It
starts $1,2,7,31,161,932,5841,38891$. The labels depend on which crossing is first and which is last,
so this is a family of meanders with distinguished ends, not closed under all symmetries of the
picture.

\subsection{The growth rate}

Let $x_0=0.9175433408\ldots$ be the positive root of $4x^4=2x+1$ and
\[ \Lambda\;=\;(1+2x_0)(1+x_0)^2\;=\;10.4245356612\ldots, \]
the larger real root of the irreducible polynomial $16x^4-112x^3-520x^2-536x+25$.

\begin{theorem}[growth rate]\label{thm:bandgrowth}
There is a constant $C_\Lambda$ such that for every $m\ge1$
\[ \Lambda^{m-1}e^{-C_\Lambda\sqrt m}\;\le\;M^{[-1,2]}_m\;\le\;\Lambda^{m-1}. \]
\end{theorem}

\begin{proof}[Key idea]
Read the crossings from left to right. After crossing $j$ the state is the current label together with
the labels of the upper arches and of the lower arches that are still open, that is, started but not
yet closed. The open arches on each side are closed in last-in-first-out order, so they form a stack.
Since the labels take only four values, each stack is a word in a finite alphabet, and the reading
process is a machine with two stacks and finitely many control states. This machine has an explicit
eigenfunction of product form, with one weight for each stack entry below the top, and
normalizing by it turns the count into a probability: $M^{[-1,2]}_m=\Lambda^{m-1}$ times the
probability that an explicit Markov chain, reading one crossing at a time, is back in its starting
state after $2m-2$ crossings without having been killed (stopped) on the way. The upper bound follows
at once. For the lower bound, consider the chain without the constraint that a stack cannot go below
empty. It has an invariant law of product form, and under it two explicit martingales show that with
probability at least $1/2$ the chain stays for $t$ steps within windows of length about $\sqrt t$ at
the tops of the two stacks. Building such windows from the start and removing them at the end costs a
factor $e^{-O(\sqrt t)}$.
\end{proof}

\subsection{The exponent}

Let $\theta=\arccos\big((1+2x_0)/(3+4x_0)\big)=1.13179\ldots$, which is $64.8468\ldots$ degrees, and
\[ \alpha_{\mathrm b}\;=\;1+\frac{\pi}{\theta}\;=\;3.7757713065\ldots . \]

\begin{theorem}[exponent]\label{thm:bandexp}
For every $\epsilon>0$ there are $c_\epsilon,C_\epsilon>0$ such that for every $m\ge2$
\[ c_\epsilon\,\Lambda^m m^{-\alpha_{\mathrm b}-\epsilon}\;\le\;M^{[-1,2]}_m\;\le\;C_\epsilon\,\Lambda^m m^{-\alpha_{\mathrm b}+\epsilon}. \]
That is, $M^{[-1,2]}_m=\Lambda^m m^{-\alpha_{\mathrm b}+o(1)}$.
\end{theorem}

Both numbers are explicit, and neither is rational: $\Lambda$ is algebraic of degree four, and
$\alpha_{\mathrm b}$ is irrational, indeed transcendental. For the irrationality, $u=2\cos\theta$ is a
root of the polynomial
\[ 113u^4-312u^3+312u^2-128u+16 , \]
obtained by substituting $x_0=(2-3u)/(4u-4)$ into $4x_0^4-2x_0-1=0$. It is irreducible: its reduction
modulo $3$, made monic, is $X^4-X-1$, which has no factor of degree one or two over the field with
three elements. So the monic minimal polynomial of $u$ has the constant term $16/113$, and $u$ is not an
\emph{algebraic integer}, a root of a monic polynomial with integer coefficients. If $\theta/\pi$ were
rational, $e^{i\theta}$ would be a root of unity and $u=e^{i\theta}+e^{-i\theta}$ would be an
algebraic integer. Hence $\theta/\pi$ is irrational, and so is $\alpha_{\mathrm b}=1+\pi/\theta$. Since
$e^{i\theta}$ is algebraic, the Gelfond--Schneider theorem then makes $\theta/\pi$, and with it
$\alpha_{\mathrm b}$, transcendental. The sharp form $M^{[-1,2]}_m\asymp\Lambda^mm^{-\alpha_{\mathrm b}}$, with
bounded ratio, is not proved, and no closed formula for the counts is known.

\subsection{The generating function}

A power series is \emph{D-finite} if it satisfies a nonzero linear differential equation whose
coefficients are polynomials. Algebraic power series are D-finite, but many D-finite series are not
algebraic. The generating function of the bounded-label family is in neither class.

\begin{corollary}\label{cor:bandnondfinite}
The generating function
\[ F_{\mathrm b}(z)\;=\;\sum_{m\ge1}M^{[-1,2]}_m z^m \]
is not D-finite, even with complex polynomial coefficients. In particular, it is not algebraic.
\end{corollary}

\begin{proof}
By the previous subsection, $\alpha_{\mathrm b}$ is irrational. Choose an integer $k$ with
$k+1>\alpha_{\mathrm b}$. The coefficient bounds of Theorem~\ref{thm:bandexp} give
\begin{equation}\label{eq:bandradial}
\lim_{r\downarrow0}\frac{\log F_{\mathrm b}^{(k)}\big((1-r)/\Lambda\big)}{\log(1/r)}\;=\;k+1-\alpha_{\mathrm b} .
\end{equation}
To see this, note that the $k$th derivative of $z^m$ is ${m(m-1)\cdots(m-k+1)}\,z^{m-k}$, and this
factor is of order $m^k$ for large $m$. So for each small $\epsilon>0$ the
positive series $F_{\mathrm b}^{(k)}((1-r)/\Lambda)$ lies, up to bounded factors and finitely many bounded
terms, between $\sum_m m^{k-\alpha_{\mathrm b}-\epsilon}(1-r)^m$ and $\sum_m m^{k-\alpha_{\mathrm b}+\epsilon}(1-r)^m$.
These sums are of order $r^{-(k+1-\alpha_{\mathrm b}\mp\epsilon)}$, which tends to infinity. Taking
logarithms and letting $\epsilon$ decrease to zero proves~\eqref{eq:bandradial}.

Suppose now that $F_{\mathrm b}$ were D-finite. Its coefficients are rational, so an equation with
complex coefficients can be replaced by one with rational coefficients: for a fixed order and degree,
the coefficients of the equation solve a system of homogeneous linear equations with rational
coefficients, and such a system with a nonzero complex solution has a nonzero rational one. Since the
coefficients of $F_{\mathrm b}$ are integers bounded by $16^m$, it is then a \emph{G-function}, a D-finite
series whose coefficients satisfy standard arithmetic growth bounds. The Andr\'e--Chudnovsky--Katz
theorem says that the minimal equation of a G-function has only regular singularities, with rational
exponents~\cite[Theorem~6]{FischlerRivoal2014}. By~\eqref{eq:bandradial}, $1/\Lambda$ is a singular point.
So near $1/\Lambda$ the function and its derivatives are finite sums of rational powers of $1/\Lambda-z$
times powers of $\log(1/\Lambda-z)$ times analytic functions. After equal powers are combined, the
first nonzero term shows that the limit in~\eqref{eq:bandradial} is rational. This contradicts the
irrationality of $\alpha_{\mathrm b}$. Finally, every algebraic power series is D-finite.
\end{proof}

The proof uses the coefficient bounds of Theorem~\ref{thm:bandexp} with arbitrary $\epsilon$; it does not
need the sharp form $M^{[-1,2]}_m\asymp\Lambda^mm^{-\alpha_{\mathrm b}}$. It also needs nothing about other
singularities on the circle $|z|=1/\Lambda$, because all estimates are made on the positive real axis.

\subsection{The wedge}

The following heuristic explains why the exponent is $1+\pi/\theta$; the proofs make it rigorous. Let $h_\uparrow(j)$ and $h_\downarrow(j)$ be the numbers of
upper and lower arches still open after crossing $j$: the heights of the two stacks. A meander starts
and ends with both heights zero, and in between neither is negative. So the pair of heights is a walk
in the quarter plane that starts and ends at the corner.

For the chain of Theorem~\ref{thm:bandgrowth}, the two heights after $t$ steps, rescaled by $\sqrt t$,
converge to a planar Brownian motion whose two coordinates have correlation
\[ \rho\;=\;-\frac{1+2x_0}{3+4x_0}\;=\;-0.42504\ldots . \]
A linear change of coordinates turns this Brownian motion into a standard one and the quarter plane
into a wedge of angle $\theta=\arccos(-\rho)$. The correlation is negative: when one stack grows, the
other tends to shrink, so it is harder for both to stay positive, and the wedge is narrower than a
right angle.

For standard Brownian motion in a wedge of angle $\theta$, started at distance one from the corner, the
probability of staying inside for time $t$ decays like $t^{-\pi/(2\theta)}$~\cite{DW}. A meander asks for
more: an excursion from the corner back to the corner. That costs $t^{-\pi/(2\theta)}$ for the first
half, the same for the second half read backward, and $t^{-1}$ for ending in a bounded region near the
corner, as for any two-dimensional walk. Altogether $t^{-1-\pi/\theta}=t^{-\alpha_{\mathrm b}}$, and
$t=2m-2$ is proportional to the order. Figure~\ref{fig:wedge} shows the picture.

\begin{figure}[t]
\centering
\begin{tikzpicture}[scale=0.7,line cap=round,line join=round]
  \fill[gray!8] (0,0)--(4.2,0) arc[start angle=0,end angle=64.846841,radius=4.2]--cycle;
  \draw[line width=1pt] (0,0)--(4.2,0);
  \draw[line width=1pt] (0,0)--(64.846841:4.2);
  \draw[-{Stealth[length=4pt]}] (3.75,0) arc[start angle=0,end angle=64.846841,radius=3.75];
  \node at (40:3.98) {$\theta$};
  \draw[blue!60!black,line width=0.6pt] plot coordinates {(0.10,0.03) (0.11,0.05) (0.14,0.05) (0.17,0.06) (0.18,0.08) (0.20,0.08) (0.21,0.10) (0.24,0.10) (0.23,0.12) (0.26,0.14) (0.28,0.13) (0.33,0.17) (0.36,0.19) (0.39,0.20) (0.38,0.19) (0.46,0.19) (0.43,0.20) (0.48,0.22) (0.50,0.27) (0.49,0.24) (0.49,0.24) (0.48,0.30) (0.51,0.32) (0.53,0.36) (0.60,0.44) (0.55,0.47) (0.57,0.51) (0.54,0.54) (0.59,0.64) (0.51,0.71) (0.50,0.71) (0.53,0.76) (0.47,0.72) (0.54,0.79) (0.57,0.84) (0.64,0.84) (0.71,0.84) (0.68,0.76) (0.67,0.82) (0.68,0.83) (0.73,0.82) (0.80,0.88) (0.78,0.86) (0.77,0.84) (0.75,0.82) (0.74,0.77) (0.75,0.91) (0.77,0.91) (0.89,0.91) (0.88,0.97) (0.89,0.83) (0.93,0.84) (0.90,0.86) (0.90,0.82) (0.92,0.86) (0.98,0.78) (0.93,0.79) (1.02,0.85) (1.02,0.75) (1.05,0.69) (1.11,0.70) (1.16,0.71) (1.23,0.69) (1.26,0.72) (1.27,0.74) (1.31,0.68) (1.27,0.69) (1.35,0.63) (1.33,0.54) (1.28,0.42) (1.40,0.48) (1.41,0.48) (1.37,0.59) (1.48,0.52) (1.48,0.51) (1.55,0.54) (1.52,0.54) (1.52,0.57) (1.45,0.58) (1.37,0.56) (1.34,0.45) (1.35,0.41) (1.40,0.50) (1.54,0.55) (1.50,0.62) (1.57,0.60) (1.55,0.60) (1.55,0.57) (1.49,0.59) (1.43,0.60) (1.41,0.64) (1.36,0.73) (1.37,0.70) (1.35,0.75) (1.42,0.68) (1.44,0.75) (1.46,0.73) (1.53,0.74) (1.56,0.75) (1.57,0.85) (1.58,0.85) (1.58,0.84) (1.64,0.79) (1.69,0.81) (1.72,0.77) (1.77,0.76) (1.75,0.76) (1.78,0.80) (1.72,0.78) (1.79,0.82) (1.81,0.91) (1.82,0.85) (1.86,0.92) (1.86,1.04) (1.86,1.02) (1.92,1.05) (1.93,0.99) (1.89,1.02) (1.92,1.01) (1.84,1.03) (1.93,0.99) (1.92,0.96) (1.93,0.94) (2.05,0.91) (2.16,0.89) (2.17,0.88) (2.20,0.83) (2.26,0.83) (2.34,0.86) (2.34,0.75) (2.35,0.76) (2.41,0.73) (2.43,0.68) (2.47,0.61) (2.50,0.57) (2.44,0.54) (2.40,0.56) (2.26,0.50) (2.25,0.55) (2.32,0.51) (2.26,0.55) (2.32,0.61) (2.37,0.55) (2.36,0.60) (2.36,0.71) (2.36,0.70) (2.32,0.73) (2.30,0.73) (2.30,0.73) (2.32,0.73) (2.36,0.71) (2.32,0.75) (2.35,0.80) (2.31,0.81) (2.30,0.89) (2.31,0.95) (2.30,0.99) (2.29,0.96) (2.30,1.06) (2.26,1.13) (2.22,1.03) (2.21,0.98) (2.25,0.96) (2.24,1.01) (2.32,1.08) (2.33,1.07) (2.17,1.09) (2.14,1.13) (2.21,1.07) (2.20,1.11) (2.17,1.10) (2.12,1.11) (2.18,1.11) (2.22,1.06) (2.24,1.01) (2.25,1.02) (2.31,1.03) (2.34,1.02) (2.27,1.08) (2.29,1.06) (2.29,1.00) (2.27,0.99) (2.30,0.98) (2.29,0.98) (2.13,1.04) (2.19,1.07) (2.23,1.12) (2.21,1.16) (2.18,1.15) (2.17,1.20) (2.10,1.19) (2.10,1.21) (2.07,1.16) (2.04,1.16) (1.98,1.15) (1.93,1.17) (1.97,1.26) (2.02,1.38) (2.03,1.43) (2.01,1.43) (1.95,1.47) (1.96,1.48) (2.03,1.45) (1.93,1.44) (1.97,1.43) (1.85,1.40) (1.81,1.47) (1.81,1.44) (1.72,1.41) (1.63,1.44) (1.62,1.41) (1.65,1.35) (1.60,1.30) (1.65,1.32) (1.59,1.34) (1.66,1.35) (1.65,1.32) (1.64,1.27) (1.65,1.31) (1.64,1.24) (1.72,1.22) (1.81,1.15) (1.81,1.13) (1.78,1.12) (1.70,1.17) (1.64,1.02) (1.62,1.09) (1.60,1.04) (1.57,1.05) (1.55,1.08) (1.55,1.04) (1.46,1.03) (1.39,1.05) (1.42,1.03) (1.44,0.99) (1.39,0.99) (1.43,0.99) (1.38,0.91) (1.34,0.97) (1.23,0.98) (1.25,0.85) (1.22,0.80) (1.25,0.88) (1.26,0.95) (1.24,0.95) (1.23,0.86) (1.27,0.84) (1.23,0.89) (1.20,0.90) (1.18,0.87) (1.20,0.83) (1.22,0.86) (1.29,0.84) (1.27,0.85) (1.31,0.86) (1.25,0.94) (1.26,0.94) (1.24,0.94) (1.30,0.84) (1.31,0.89) (1.27,0.92) (1.27,0.91) (1.27,0.79) (1.23,0.75) (1.19,0.74) (1.12,0.84) (1.23,0.86) (1.22,0.88) (1.24,0.91) (1.16,0.97) (1.09,0.90) (1.06,0.85) (1.03,0.78) (1.04,0.75) (0.97,0.72) (0.97,0.71) (1.03,0.65) (0.96,0.64) (0.98,0.59) (0.91,0.68) (0.91,0.64) (0.90,0.76) (0.96,0.83) (0.91,0.75) (0.89,0.79) (0.92,0.80) (0.84,0.71) (0.82,0.73) (0.75,0.71) (0.69,0.71) (0.70,0.66) (0.71,0.64) (0.80,0.64) (0.82,0.61) (0.84,0.57) (0.80,0.55) (0.76,0.49) (0.73,0.49) (0.62,0.38) (0.65,0.42) (0.60,0.41) (0.53,0.41) (0.49,0.41) (0.49,0.36) (0.49,0.40) (0.45,0.36) (0.44,0.35) (0.41,0.33) (0.36,0.34) (0.33,0.32) (0.31,0.27) (0.26,0.26) (0.25,0.20) (0.20,0.18) (0.18,0.16) (0.15,0.16) (0.12,0.13) (0.10,0.11) (0.09,0.09) (0.07,0.07)};
  \fill (0,0) circle[radius=1.4pt];
  \node[below left,font=\small] at (0,0) {corner};

\end{tikzpicture}
\caption{Schematic: after a linear change of coordinates the two stack heights move like a planar
Brownian motion in a wedge of angle $\theta$, and a meander is an excursion from the corner back to the
corner. Here $\theta\approx64.85^\circ$. The path drawn is illustrative, not a sample.}
\label{fig:wedge}
\end{figure}

Two checks make the formula plausible. Without the requirement of a single loop, the upper and lower
arch systems are independent, the heights are independent, $\rho=0$, $\theta=\pi/2$, and the exponent
is $1+2=3$; indeed the number of pairs of arch systems is $\mathrm{Cat}_n^2\approx16^n/(\pi n^3)$. For the
narrower band of labels $\{-1,0,1\}$, the counts $1,2,6,22,92,422,2074,\ldots$ are the Baxter
numbers. Chung, Graham, Hoggatt and Kleiman gave an explicit formula for them~\cite{CGHK}, from which
they grow like a constant times $8^n n^{-4}$;
there $\rho=-1/2$, $\theta=\pi/3$ and $1+\pi/\theta=4$. The identification with the Baxter numbers is
from the literature and was checked here only on small orders, not proved.

The main difficulty in the proofs is that the heights are not a random walk with independent steps:
the chain is Markov only on its full state, two words of unbounded length. The proofs replace
independence by stationarity under the invariant law, and the lower bound also uses that the chain
reversed in time is again an explicit two-stack machine.

For all meanders the stacks are unbounded in both directions and no finite description of this kind
is known. Whether the exponent of all meanders arises from a wedge in a similar way is the question
behind Papers~A and~D~\cite{PaperA,PaperD}.
